\documentclass[11pt]{article}

\input{0-parameters.tex}

\begin{document}
\section{Appendix: Switching from Gas to Renewable Electricity (Final, Useful, and Primary)}
\label{appendix_gas_renewables}

Energy accounting typically distinguishes three stages of energy flows: \textbf{primary}, \textbf{final}, and \textbf{useful}. \textbf{Primary energy} refers to raw energy sources (e.g., coal, crude oil, wind) \emph{before} any transformation. \textbf{Final energy} is the energy commodity \emph{delivered} to end-use sectors after conversion and transmission. \textbf{Useful energy} is the \emph{portion} of final energy that provides actual energy services (heat, mechanical work, lighting), accounting for end-use inefficiencies \citep{wayEmpiricallyGroundedTechnology2022}.

\vspace{3mm}
This appendix serves as an exercise in verifying accounting identities using simplified examples from the energy sector, under assumptions such as a uniform price for useful energy across all fuel types. It examines how the greening of electricity production (transition from gas-based to renewable electricity, Section~\ref{sect_greening_elec}) and the electrification of manufacturing energy use (Section~\ref{sect_energy_electrification}) affect sectoral and aggregate productivity measures.  
\textbf{We assess whether the choice of energy accounting basis—primary, final, or useful—affects productivity accounting in these two cases.}

\vspace{3mm}

We compare productivity growth in two economic structures, using the three conventions of energy accounting, and two productivity accounting frameworks.
First, an economy with three separate energy sectors: $G$ (natural gas supply), $F$ (fossil-based electricity production), and $R$ (renewable electricity supply) (Section~\ref{sect_energy3sectors}). Second, an economy with a single aggregated energy sector supplying gas and both types of electricity (Sections~\ref{sect_greening_elec} and \ref{sect_energy_electrification}). In both cases, the manufacturing sector consumes natural gas and electricity.

We compare productivity changes using two frameworks: the primal--dual TFP system (Eqs.~\eqref{eq:primal_appendix}--\eqref{eq:dual_appendix}) and the $\zeta$-framework (Eq.~\eqref{eq:zeta_framework}):

\begin{subequations}
\begin{align}
\text{Primal: }\hat{A}_i  &=  
\hat{X}_i  -  \sum_{j=1}^N a_{ji} \hat{X}_{ij}  -  \sum_{f=1}^F \tilde{l}_{fi} \hat{L}_{fi}
\label{eq:primal_appendix} \\
\text{Dual: }\hat{A}_i  &= 
-\left( \hat{p}_i  -  \sum_{j=1}^N a_{ji} \hat{p}_j  -  \sum_{f=1}^F \tilde{l}_{fi} \hat{w}_f \right)
\label{eq:dual_appendix} \\
\zeta\text{-framework: }\hat{\tilde{A}}_i &= 
\zeta \hat{X}_i+(\zeta-1)\hat{p}_i
\label{eq:zeta_framework}
\end{align}
\end{subequations}


\paragraph{Unit conversion} Let $\beta$ denote the efficiency of converting gas to electricity, and $\alpha$ the efficiency from final electricity to useful energy, with $0 \leq \alpha, \beta \leq 1$. One MWh of natural gas produces $\beta$ MWh of final electricity and $\alpha\beta$ MWh of useful energy, while one MWh of renewable electricity corresponds directly to one MWh of both primary and final energy, delivering $\alpha$ MWh of useful energy.
\[\text{gas}~ \xrightarrow{\beta}~ \text{electricity}~ \xrightarrow{\alpha} ~\text{useful}\]

\paragraph{Price conversions}
We assume electricity prices are equal across sources, fossil $F$ and renewables $R$, $p_R = p_F$. Producing 1~kWh of fossil electricity requires $1/\beta$~kWh of natural gas, at unit cost $p_G$, so delivering $X_F$ units  of electricity, at unit cost $p_F$, costs $p_G \frac{X_F}{\beta} = p_F X_F$, implying $p_F = \frac{\text{cost}}{\text{quantity}}=\frac{p_G \sfrac{X_F}{\beta}}{X_F} =\sfrac{p_G}{\beta}$. Similarly, the unit cost of useful energy is $p_R/\alpha$. Throughout the appendix, we assume gas and renewable prices are constant, $\hat{p}_G = \hat{p}_R = 0$.


\subsection{Summary of Results}



First, when energy sectors consistently supply distinct energy types (gas or electricity), productivity remains unchanged under greening (renewables substitution) or electrification.
Second, when considering a single energy sector and a switch from natural gas to renewable electricity, national accounting correctly shows zero productivity growth at all stages. By contrast, under the final and primary energy approaches, applying the $\zeta$-framework introduces a bias at either the electrification or greening step. Only the useful energy approach ensures full consistency throughout the process.\footnote{
In a single aggregated energy sector, replacing gas with electricity reduces total sectoral output but delivers the same useful energy services. In contrast to the setting discussed in Appendix~(A7), where changes in an upstream sector affect downstream sectors through output linkages, the energy sector here provides constant useful energy, so a reduction in inputs does not imply lower downstream output.
} 


\subsection{Distinct energy sectors}
\label{sect_energy3sectors}
We consider three separate energy sectors: gas, gas-based electricity, renewables-based electricity. We show that productivity is unaffected, irrespective of the approach to energy -- final, primary or useful -- when demand is transferred from one sector to another through  electrification ($G\rightarrow F$), greening electricity ($F\rightarrow R$), or rapid decarbonisation ($G\rightarrow R$).

We assume that for each sector $i\in\{G,F,R\}$, the output shrinks by the factor $\delta_i$,
$\hat{X}_i= \ln\bigl((1-\delta_i) X_i\bigr)  -  \ln(X_i)= \ln(1-\delta_i)$.
We assume all inputs to $i$ scale down proportionally to the output, so each input grows at a rate $\ln(1-\delta_i)$. From Eq.~\ref{eq:primal_appendix}, hence productivity growth is 
\[
    \hat{A}_i
= \ln(1-\delta_i)
  - \sum_{j=1}^N a_{ji}\ln(1-\delta)
  - \sum_{f=1}^F \tilde{l}_{fi}\ln(1-\delta)
= \ln(1-\delta_i)  -  \ln(1-\delta_i)
= 0.
\]

From the dual approach (Eq.~\ref{eq:dual_appendix}), $\hat{A}_i=-\hat{p}_i-\hat{\Sigma}_i$ and null productivity growth, we get that $\hat{p}_i=0$, assuming the price of all inputs remains constant. 

\vspace{3mm}
The manufacturing sector faces an unchanged energy bill when part of its energy demand shifts from one source to another. In the primary energy approach, a shift between $F$ and $R$, that is shifting a quantity $\eta$ of electricity consumption from gas-based to renewable sources corresponds to replacing $\eta/\beta$ units of gas (priced at $p_G$) with $\eta$ units of renewable electricity (priced at $p_R = p_G/\beta$). The higher unit cost of renewable electricity exactly offsets the higher volume of gas required, leaving total spending unchanged. Likewise, when shifting a quantity $\eta/\beta$ from Gas to renewable electricity.

In the final and useful energy approaches, transferring a quantity $\eta$ from gas ($G$) to renewable ($R$) or fossil electricity ($F$) leads to an increase in electricity production by $\eta\beta$. The cost shifts from $\eta p_G$ to $\eta\beta p_R = \eta\beta (p_G/\beta) = \eta p_G$, again leaving expenditures unchanged. A shift between $F$ and $R$ in final or useful energy is similarly iso-cost.

We now detail the impact on final, primary, and useful energy quantities. Since the manufacturing sector's output $\hat{X}_M$ remains constant and its energy bill is unchanged ($\hat{S}_M = 0$), its price also remains unchanged ($\hat{p}_M = 0$). With input prices constant, productivity growth is zero by the dual approach (Eq.~\eqref{eq:dual_appendix}). As all sectors exhibit zero productivity growth, aggregate productivity growth is likewise zero.

\vspace{3mm}

In our framework, sectoral productivity changes as $\hat{\tilde{A}}_i = \zeta\hat{X}_i=\zeta \ln(1-\delta_i)$. We establish the following identity:
\[
\sum_i p_i X_i \ln(1-\delta_i) = 0,
\]
which can be derived as follows, noting that $\hat{S}_i=\hat{X}_i+\hat{p}_i=\hat{X}_i=\ln(1-\delta_i) $
\begin{equation}
\label{eq_weighted_share_change_cost}
\begin{aligned}
\sum_i p_i X_i \ln(1-\delta_i) 
&= \sum_i S_i \hat{S}_i \\
&= \sum_j S_j\times  \sum_i \frac{S_i}{\sum_j S_j} \hat{S}_i \\
&= \sum_j S_j\times  \widehat{\left(\sum_i S_i\right)} = 0,
\end{aligned}
\end{equation}
because the total cost $\sum_i S_i$ remains unchanged before and after the switch from one source to another.\footnote{
Recall that for two variables $x$ and $y$, the log change satisfies $\widehat{x+y} = \frac{x}{x+y} \hat{x} + \frac{y}{x+y} \hat{y}$.
}

\noindent Hence, denoting $pY$ as GDP, aggregate productivity growth is for all $\zeta$
\[
\hat{\tilde{A}} = \lambda_M \hat{\tilde{A}}_M + \lambda_F \hat{\tilde{A}}_F + \lambda_R \hat{\tilde{A}}_R + \lambda_G \hat{\tilde{A}}_G = \frac{\zeta}{pY} \sum_{i\in\{G,F,R\}} p_i X_i \ln(1-\delta_i) = 0.
\]

\vspace{3mm}

Therefore, regardless of the approach — primary, useful, or final — as long as different energy sub-sectors are consistently defined in terms of outputs, inputs, and prices, productivity remains unaffected by either the greening of electricity (shifting from gas-based to renewable electricity) or electrification (shifting from gas to electricity, whether renewable or gas-based).


\subsection{Greening electricity}
\label{sect_greening_elec}
We consider in this section and the next, a single energy sector providing energy in the form of natural gas ($X_G$) and electricity, from fossil fuel ($X_F$) and renewables ($X_R$). We examine the shift a quantity $\eta$ of fossil electricity ($F$) to renewables electricity ($R$). We distinguish the useful energy conversion coefficient for electricity $\alpha$ and natural gas, $\alpha_G$.

Total production of energy $X_E^f,X_E^p,X_E^u$ expressed respectively in final, primary and useful units are initially
\[
X_E^f=X_G+\beta X_F+X_R, \quad X_E^p=X_G+ X_F+X_R, \quad \text{and} \quad  X_E^u=\alpha_G X_G+ \alpha\beta X_F+\alpha X_R.
\]
After the switch of $\eta$ units from $G$ to $F$, production is  $X_E^{f'},X_E^{p'},X_E^{u'}$, showing that $\hat{X}_E^f=\hat{X}_E^u=0$, but that $\hat{X}_E^p<0$,
\begin{equation*}
    \begin{cases}
        \text{Final: }& X_E^{f'}=X_G+\beta X_F+X_R + (-\eta + \eta)=X_E^f,\\
         \text{Primary: }& X_E^{p'}= X_G+ X_F+X_R +  \eta(1-\frac{1}{\beta}), \\
        \text{Useful: }& X_E^{u'}=\alpha_G X_G+ \alpha X_F+\alpha X_R + \alpha(-\eta + \eta )= X_E^u.  \\
    \end{cases}
\end{equation*}


Energy sales are equal in the three approaches (see Section \ref{sect_energy3sectors}),
$S_E=p_GX_G+ X_Fp_F+X_Rp_R$. Using $\hat{S}_E = \hat{X}_E + \hat{p}_E$, we get price growth, $\hat{p}_E^f = 0$, $\hat{p}_E^u = 0$, and $\hat{p}_E^p = -\hat{X}_E^p, which is positive.$

\paragraph{Productivity Accounting.} Using the national accounting approach, from Eq.~\ref{eq:primal_appendix},  we assume all inputs to each of the three energies $i$ $\in \{G,F,R\}$, scale down proportionally to the output at a rate $1-\delta_i$. Since, the total cost to produce each of the three sub-products also change proportionally to the output (see previous section \ref{sect_energy3sectors}), we can show,\footnote{
We separate the inputs of the energy sector depending on whether they are used to produce the energy $i$, gas, or electricity from gas or renewables. We denote $a_{ji}$ the ratio of the cost of  the input $j$ to produce the sub-product $i$ and the total cost of the energy sector, likewise for $\tilde{l}_{fi}$. We sum up the change in all inputs quantities weighted by their cost share as in Eq.~\ref{eq:primal_appendix},
\begin{equation*}
\begin{aligned}
\sum_{i\in \{G,F,R\}}\left( \sum_{j=1}^N a_{ji} \hat{X}_{ij}  +  \sum_{f=1}^F \tilde{l}_{fi} \hat{L}_{fi}\right)&=
\sum_{i\in \{G,F,R\}}\ln(1-\delta_i)\left( \sum_{j=1}^N a_{ji}   +  \sum_{f=1}^F \tilde{l}_{fi}\right)\\&=
 \sum_{i\in \{G,F,R\}}\hat{S}_i\left( \frac{S_i}{\sum_{i\in \{G,F,R\}} S_i }\right)=
\widehat{\left(\sum_{i\in \{G,F,R\}}\hat{S}_i\right)}= 
 0.
\end{aligned}
\end{equation*}
} 
as in Eq. \eqref{eq_weighted_share_change_cost}, that the cost-weighted input growth, $\hat{\Sigma}_X$, of the whole energy sector $E$ is 
\begin{equation}
\label{eq_weighted_share_PGR}
   \hat{\Sigma}_X= \left(\sum_{j=1}^N a_{jE} \hat{X}_{Ej}  +  \sum_{f=1}^F \tilde{l}_{fE} \hat{L}_{fE}\right)=0.
\end{equation}

Thus, $\hat{A}_E=\hat{X}_E$. The energy sector productivity growth, under the three metrics, are $\hat{A}_E^f = 0$, $\hat{A}_E^u = 0$, and $\hat{A}_E^p =\hat{X}_E^p=-\hat{p}_E^p$.

In the manufacturing sector, since the final and useful energy demand and prices remain unchanged, productivity is unaffected, that is, $\hat{A}_M^f = \hat{A}_M^u = 0$. When considering primary energy, however, the quantity of energy decreases. From the primal (or equivalently, dual) definition of productivity (Eq.~\ref{eq:primal_appendix}), we have $\hat{A}_M = \hat{X}_M - a_{EM} \hat{X}_E^p = -a_{EM} \hat{X}_E^p$.

In the final and useful energy approaches, productivity remains unchanged in both the manufacturing and energy sectors, and thus aggregate productivity is also unchanged. Under primary energy accounting, the productivity decline in the energy sector is exactly offset by the productivity gain in the manufacturing sector, so that
\[
\hat{A}^p=\lambda_E \hat{A}_E^p+\lambda_M \hat{A}_M^p=\frac{X_Ep_E}{pY} \hat{X}_E^p+ \frac{p_XX_M}{pY} (-a_{EM}\hat{X}_E^p)=\frac{X_Ep_E}{pY} \hat{X}_E^p - \frac{p_XX_M}{pY} \frac{X_Ep_E}{p_XX_M}\hat{X}_E^p=0.
\]
Note that the Domar weights are not specific to the approach - primary, final or useful - as the sales and GDP remain unchanged. 

\paragraph{Framework Derivation.}
In our framework, the productivity of the energy sector is $\hat{\tilde{A}}_E^f = 0$, and $\hat{\tilde{A}}_E^u = 0$, and  $\hat{\tilde{A}}_E^p=\zeta \hat{X}_E^p+(\zeta-1) \hat{p}_E^p=\zeta \hat{X}_E^p+(1-\zeta) \hat{X}_E^p=\hat{X}_E^p$. In all approaches, the manufacturing sector pays the same energy bill ($\hat{S}_E^p$), we assume $\hat{X}_M=0$, then $\hat{p}_M=\hat{S}_M-\hat{X}_M= 0$. Hence, $\hat{\tilde{A}}_M=0$. Aggregate productivity is then $\hat{\tilde{A}}^f=0$, $\hat{\tilde{A}}^u=0$, and $\hat{\tilde{A}}^p = \lambda_E \hat{X}_E^p$.

\subsection{Electrification}
\label{sect_energy_electrification}

As in Section~\ref{sect_greening_elec}, we consider a single energy sector supplying natural gas ($X_G$) and electricity, produced from fossil fuels ($X_F$) and renewables ($X_R$). We examine the shift of a quantity $\gamma$ of gas ($G$) to fossil electricity ($F$), corresponding to $\gamma\beta$ units of electricity.  

In the useful energy accounting, this shift leaves the aggregate energy quantity unchanged. Similarly, under primary energy accounting, the same quantity of gas $\gamma$ is used to deliver $\gamma\beta$ of electricity, so there is no net change of productivity. However, differences arise under the final energy approach.  

\vspace{3mm}

In the case of electrification, we get $\hat{X}_E^p=\hat{X}_E^u=0$, but $\hat{X}_E^f<0$, as
\begin{equation*}
    \begin{cases}
        \text{Final: }& X_E^{f'}=X_G+\beta X_F+X_R + \gamma(\beta-1),\\
         \text{Primary: }& X_E^{p'}= X_G+ X_F+X_R +  (-\gamma+\gamma)=X_E^p, \\
        \text{Useful: }& X_E^{u'}=\alpha_G X_G+ \alpha X_F+\alpha X_R + \alpha(-\gamma + \gamma )= X_E^u.  \\
    \end{cases}
\end{equation*}
Again, the sales are equal in the three approaches and do not change with the electrification. Hence, $\hat{S}_E^f=\hat{S}_E^u=\hat{S}_E^p=0$. We get price growth from $\hat{p}_E=\hat{S}_E - \hat{X}_E $, $\hat{p}_E^f = -\hat{X}_E^f$, 
 $\hat{p}_E^u = 0$, and $\hat{p}_E^p = 0$.

\paragraph{Productivity Accounting.} Using the national accounting approach, from Eq.~\ref{eq:primal_appendix},  we assume all inputs to each of the three energies $i$ $\in \{G,F,R\}$, scale down proportionally to the output at a rate $1-\delta_i$ and thus show that the cost-weighted input growth $\hat{\Sigma}_X$ of the whole energy sector $E$ is null, as in Eq. \eqref{eq_weighted_share_PGR}. Thus, the energy sector productivity growth, under the three metrics, are $\hat{A}_E^f = \hat{X}_E^f$, $\hat{A}_E^u = 0$, and $\hat{A}_E^p = 0$. Likewise, for the manufacturing sector $
\hat{A}_M^f = \hat{X}_M-a_{EM}\hat{X}_E^f= -a_{EM}\hat{X}_E^f$, $\hat{A}_M^p = 0$, $\hat{A}_M^u = 0$.
Once again, the manufacturing and energy sector productivity changes offset one another and aggregate productivity growth is null, $\hat{A}^f=\lambda_E \hat{A}_E^f+\lambda_M \hat{A}_M^f=0$.

\paragraph{Framework Derivation.}

Our framework-based energy productivity gives $\hat{\tilde{A}}_E^f = 0$, $\hat{\tilde{A}}_E^u = 0$, and  $\hat{\tilde{A}}_E^f=\zeta \hat{X}_E^f+(\zeta-1) \hat{p}_E^f=\zeta \hat{X}_E^f+(1-\zeta) \hat{X}_E^f=\hat{X}_E^f$.
In all approaches, the manufacturing sector access the same energy service, hence $\hat{\tilde{A}}_M=0$. Aggregate productivity is then, $\hat{\tilde{A}}^f=\lambda_E \hat{X}_E^f$, $\hat{\tilde{A}}^u=0$, and 
$\hat{\tilde{A}}^p = 0$.


% Bibliography ---------------------------

\tiny
\renewcommand\bibname{References}
\bibliographystyle{apalike}
\bibliography{z-bib}
\normalsize
\end{document}


% \subsection{Why is it different from previous appendix? (months ago)}
% \ER{Subsection to be deleted after François' reading}

% NB: in the previous appendix, why did I find it didn't change anything ? Because I was using our framework approach for the energy sector (with the subcase of a full switch, thus $\delta=1$) 
% \[ \hat{A}_E=\hat{X}_E=\ln(\delta \beta )=\ln(\beta ),\]
% However, for the manufacturing sector I was using the true productivity measure, 

%   \begin{equation*}
%     \left\{
%     \begin{aligned}
%         \hat{\Sigma}_X &= w_{EM}\hat{X}_{ME} 
%                       = w_{EM}\ln \beta , \\
%         \hat{\Sigma}_p &= w_{EM}\hat{p}_E
%                       = w_{EM}\bigl(-\ln \beta \bigr).
%     \end{aligned}
%     \right.
% \end{equation*}

% where $w_{EM}$ is the energy share in M's cost.  Assuming $\hat{X}_M = 0$, one derives $\hat{p}_M = 0$ and
% \begin{equation*}
%     \hat{A}_M = -w_{EC}\ln \beta .
% \end{equation*}
% And then, in the aggregate, it cancels out. 
% \[\hat{A}=0\]

% But if we use the framework approach as well for the manufacturing sector, then because $\hat{S}_M=0$\footnote{Because$w_{EC}$  remains constant, as the manufacturing sector's total energy bill does not change: 
% \begin{equation*}
%      p_GX_F 
%          = \frac{p_R}{\beta }\beta X_F 
%          = p_RX_R.
% \end{equation*} Hence E's share in M's costs stays the same before and after.
% }

% \[
% \left\{
% \begin{aligned}
% \hat{X}_M&=0\\
% \hat{p}_M &= \hat{S}_M-\hat{X}_M= 0\\
% \hat{A}&=0\\
% \end{aligned}
% \right.
% \]

% \newpage



